\documentclass[../main.tex]{subfiles}
\begin{document}

\section*{Kết quả đạt được}
Dự án đã hoàn thành một hệ thống quản lý dự án chạy được và đáp ứng các mục tiêu đề ra:
\begin{itemize}
  \item Tài khoản không có vai trò toàn cục; dự án có vai trò PM và Thành viên theo từng dự án, lời
        mời phải được chấp nhận, và server kiểm tra quyền ở mọi route.
  \item Bảng Kanban và dạng bảng ưu tiên bàn phím, sắp xếp theo hạn chót, nhiều người thực hiện mỗi
        task, phụ thuộc, đường găng và màn hình DAG.
  \item Quy trình Scrum: vòng đời sprint, lập kế hoạch backlog, burndown, velocity và retrospective
        do AI viết; cùng lịch sử thay đổi cho mọi task.
  \item Năm chức năng AI dùng Groq với dự phòng ba tầng, đầu ra có cấu trúc được kiểm tra, và trả lời
        theo ngôn ngữ giao diện.
  \item Hoạt động local-first với chính sách ngoại tuyến an toàn; giao diện Anh/Việt hoàn chỉnh.
  \item Bản triển khai production tại \url{https://koshi.huynhca.com}, 185 test backend và 442 test
        frontend, tích hợp liên tục, và tám tài liệu sống truy vết mọi yêu cầu tới mã và test.
  \item Một ``nhà máy phần mềm'' agent AI: sáu trạm agent chuyên trách, rào chắn thực thi chính sách
        dự án trên mọi lần sửa và mọi lệnh (56 test), cùng health check, mã yêu cầu và metrics để vận
        hành sản phẩm.
\end{itemize}

Ngoài các chức năng, bài học giá trị nhất là về kiểm chứng. Nhiều lỗi nghiêm trọng nhất---khoá bí
mật nằm trong image container, lệnh triển khai suýt ghi đè production bằng cấu hình phát triển, một
chức năng ``AI'' chưa từng gọi mô hình, nút chuyển ngôn ngữ để lại nửa giao diện chưa dịch---đều vô
hình trước một bộ test toàn xanh. Chúng chỉ được tìm thấy khi kiểm tra sản phẩm thật: soi image đã
build, đọc lệnh trước khi chạy, gọi mô hình thật, quét từng màn hình. Thói quen kiểm thử đột biến cho
chính các test, và ghi lại mọi quyết định cùng lý do, bắt nguồn từ những trải nghiệm đó.

\section*{Hạn chế}
\begin{itemize}
  \item Dự án chung chỉ đọc khi ngoại tuyến; chưa có thuật toán gộp các chỉnh sửa đồng thời ngoại
        tuyến.
  \item Burndown được tính theo danh sách task hiện tại của sprint, nên việc bị chuyển ra khỏi sprint
        sẽ biến mất khỏi lịch sử của sprint đó.
  \item Câu trả lời dự phòng tất định của AI hiện chỉ có tiếng Việt.
  \item Bản công khai chưa giới hạn số lần đăng nhập sai, và các chỉ tiêu hiệu năng trong yêu cầu
        (hiển thị trong một khung hình, dung lượng bộ nhớ) chưa được đo.
  \item Bản công khai chạy trên một máy duy nhất và không truy cập được khi máy đó tắt.
\end{itemize}

\section*{Hướng phát triển}
\begin{itemize}
  \item Dùng CRDT (kiểu dữ liệu nhân bản không xung đột) để dự án chung cũng sửa được khi ngoại tuyến.
  \item Nhật ký sự kiện theo sprint để có lịch sử burndown chính xác, và màn hình daily scrum.
  \item Giới hạn tần suất đăng nhập, khoá tài khoản; chuyển sang hosting có sao lưu tự động.
  \item Ứng dụng di động gốc và tích hợp với nền tảng Git (tự động liên kết commit, pull request với
        task).
\end{itemize}

\end{document}
